\begin{enumerate}

\item Velocidad periférica

\[\fitmath{U=\Omega\cdot \mathrm{rad}}\]

\[\fitmath{U=(1010)\cdot (0.271)\approx 273.71\,\mathrm{m/s}}\]

\item La velocidad tangencial absoluta sigue vórtice libre.

\[\fitmath{\mathrm{ct}=\frac{\mathrm{kv}}{\mathrm{rad}}}\]

\[\fitmath{\mathrm{ct}=\frac{(14)}{(0.271)}\approx 51.66052\,\mathrm{m/s}}\]

\item Se forma el ángulo del gas en referencia rotor.

\[\fitmath{\beta=\arctan\left(\frac{\frac{\mathrm{kv}}{\mathrm{rad}}-\Omega\cdot \mathrm{rad}}{\mathrm{cx}}\right)}\]

\[\fitmath{\beta=\arctan\left(\frac{\frac{(14)}{(0.271)}-(1010)\cdot (0.271)}{(159)}\right)\approx -0.9493752\,\mathrm{rad}}\]

\item Se despeja metal a partir de incidencia definida.

\[\fitmath{\mathrm{kap}=\arctan\left(\frac{\frac{\mathrm{kv}}{\mathrm{rad}}-\Omega\cdot \mathrm{rad}}{\mathrm{cx}}\right)-\mathrm{inc}}\]

\[\fitmath{\mathrm{kap}=\arctan\left(\frac{\frac{(14)}{(0.271)}-(1010)\cdot (0.271)}{(159)}\right)-(-0.01745329252)\approx -0.9319219\,\mathrm{rad}}\]

\end{enumerate}
